%% paper2.tex — Development System: архитектура многопотоковой LLM-системы
%% Компиляция: lualatex paper2.tex
\documentclass[11pt,a4paper]{article}

\usepackage{fontspec}
\setmainfont{DejaVu Serif}
\setmonofont[Scale=0.88]{DejaVu Sans Mono}

\usepackage{geometry}
\geometry{
  a4paper,
  top=25mm, bottom=25mm,
  left=28mm, right=28mm,
  headsep=10mm
}

\usepackage{polyglossia}
\setmainlanguage{russian}
\setotherlanguage{english}

\usepackage{microtype}
\usepackage{parskip}
\usepackage{xcolor}
\usepackage{booktabs}
\usepackage{enumitem}
\usepackage{titlesec}
\usepackage{fancyhdr}
\usepackage{hyperref}
\usepackage{caption}
\usepackage{float}
\usepackage{amsmath}

% ─── листинги ────────────────────────────────────────────────────────────────
\usepackage{listings}

\definecolor{clbg}{HTML}{F8F8F2}
\definecolor{clcomment}{HTML}{6A9955}
\definecolor{clkeyword}{HTML}{569CD6}
\definecolor{clstring}{HTML}{CE9178}
\definecolor{clparen}{HTML}{999999}
\definecolor{clbuiltin}{HTML}{DCDCAA}
\definecolor{clpackage}{HTML}{4EC9B0}

\lstdefinelanguage{CommonLisp}{
  morekeywords={defun,defpackage,defvar,defparameter,defmacro,defstruct,
                let,let*,labels,flet,lambda,cond,when,unless,case,
                progn,handler-case,handler-bind,restart-case,ignore-errors,
                loop,do,dolist,dotimes,mapcar,mapc,reduce,remove-if,
                format,write,read,load,require,
                in-package,use-package,export,import,
                setf,setq,push,pop,incf,decf,
                values,multiple-value-bind,
                and,or,not,null,if,t,nil,
                car,cdr,cons,list,append,reverse,length,
                with-open-file,with-output-to-string,
                error,warn,signal,funcall,apply,
                string=,string-upcase,find-symbol,find-package,fboundp,
                pathname,merge-pathnames,probe-file,ensure-directories-exist,
                sb-ext:exit,sb-ext:posix-getenv,uiop:getenv,uiop:run-program,
                asdf:defsystem,asdf:load-system,ql:quickload},
  sensitive=false,
  morecomment=[l]{;},
  morecomment=[s]{\#|}{|\#},
  morestring=[b]",
}

\lstset{
  language=CommonLisp,
  backgroundcolor=\color{clbg},
  basicstyle=\footnotesize\ttfamily,
  keywordstyle=\color{clkeyword}\bfseries,
  commentstyle=\color{clcomment}\itshape,
  stringstyle=\color{clstring},
  frame=single,
  framesep=4pt,
  rulecolor=\color{gray!30},
  showstringspaces=false,
  breaklines=true,
  breakatwhitespace=true,
  tabsize=2,
  numbers=left,
  numberstyle=\tiny\color{gray},
  numbersep=6pt,
  xleftmargin=18pt,
  xrightmargin=4pt,
  captionpos=b,
  literate=
    {←}{{$\leftarrow$}}1
    {→}{{$\rightarrow$}}1,
}

\lstnewenvironment{lispcode}[1][]
  {\lstset{language=CommonLisp,#1}}
  {}

% ─── оформление ──────────────────────────────────────────────────────────────
\hypersetup{
  colorlinks=true,
  linkcolor=black,
  urlcolor=blue!60!black,
  citecolor=black
}

\titleformat{\section}{\large\bfseries}{\thesection.}{0.5em}{}
\titleformat{\subsection}{\normalsize\bfseries}{\thesubsection.}{0.5em}{}
\titleformat{\subsubsection}{\normalsize\itshape}{\thesubsubsection.}{0.5em}{}

\pagestyle{fancy}
\fancyhf{}
\fancyhead[L]{\small\textit{Development System: архитектура LLM-системы на Common Lisp}}
\fancyhead[R]{\small\thepage}
\renewcommand{\headrulewidth}{0.4pt}

% ─── документ ────────────────────────────────────────────────────────────────
\begin{document}

\begin{titlepage}
\centering
\vspace*{2cm}
{\LARGE\bfseries Development System\\[0.4em]
\large Архитектура многопотоковой LLM-системы\\
с персистентным образом на Common Lisp}

\vspace{1.5cm}
{\large Артемий Уродовских}\\[0.3em]
{\normalsize 2026}

\vspace{2cm}
\begin{abstract}
\noindent
Статья описывает \textbf{Development System} — двухкомпонентную систему для оркестрации вычислительных потоков с языковой моделью в роли исполнителя.
Система состоит из \textit{Мастера} — Telegram-бота, написанного на Common Lisp и принимающего команды, — и \textit{Подмастерья} — персистентного SBCL-образа, в который потоки загружаются динамически без перезапуска процесса.

Ключевая идея: поток — это автономный модуль Common Lisp, определяющий единственную функцию \texttt{выполнить/1}. Мастер знает, как потоки запускать и комбинировать; Подмастерье знает, как они загружаются в живой образ и вызываются напрямую, минуя subprocess-оверхед.

Особого внимания заслуживает \textit{поток-порождатель} — поток, который с помощью LLM генерирует новые потоки, компилирует их и загружает в тот же образ, тем самым замыкая метациркулярную петлю исполнения.

Система полностью мигрирована с Racket на Common Lisp. Общий объём кодовой базы — около 1100 строк, система проходит 6 smoke-тестов.
\end{abstract}

\vspace{1cm}
\noindent\textbf{Ключевые слова:} Common Lisp, SBCL, языковые модели, Telegram-бот, динамическая загрузка кода, персистентный образ, метациркулярность
\end{titlepage}

\tableofcontents
\newpage

% ═══════════════════════════════════════════════════════════════════════════════
\section{Введение}

Системы, управляемые языковыми моделями, обычно строятся вокруг одного из двух паттернов: либо LLM вызывается синхронно как функция («ответь на вопрос»), либо запускается автономный агент с инструментами. Оба подхода предполагают, что вычислительная единица — это HTTP-вызов к API.

\textit{Development System} предлагает иную структуру: \textbf{поток} (\textit{stream}) — это скомпилированный модуль Common Lisp, работающий в одном долгоживущем процессе SBCL. Взаимодействие с пользователем происходит через Telegram-бот. Внутри — оркестратор и исполнитель, разделённые по ответственности, но работающие в одном языке.

Переход с Racket на Common Lisp продиктован прагматическими соображениями:
\begin{itemize}[noitemsep]
  \item SBCL поддерживает горячую замену кода (\texttt{load}) в живом процессе без перезапуска;
  \item Common Lisp предоставляет мощную систему условий (\textit{condition system}) для структурированной обработки ошибок;
  \item Единый язык в стеке снижает когнитивную нагрузку: не нужно переключаться между Racket-семантикой в оркестраторе и CL-семантикой в исполнителе;
  \item Racket-процесс нестабильно завершался при зависших HTTP-соединениях — \texttt{user break} прерывал poll-loop.
\end{itemize}

\subsection{Структура статьи}

Раздел~\ref{sec:arch} описывает общую архитектуру. Раздел~\ref{sec:master} — Мастер: Telegram API, роутинг команд, управление состоянием. Раздел~\ref{sec:apprentice} — Подмастерье: персистентный образ и загрузка потоков. Раздел~\ref{sec:spawner} — метациркулярный порождатель. Раздел~\ref{sec:tracer} — механизм трассировки. Раздел~\ref{sec:design} разбирает ключевые архитектурные решения. Раздел~\ref{sec:conclusion} — заключение.

% ═══════════════════════════════════════════════════════════════════════════════
\section{Архитектура системы}
\label{sec:arch}

\subsection{Компоненты}

Система состоит из двух SBCL-процессов и набора файлов-потоков (рис.~\ref{fig:arch}).

\begin{figure}[H]
\centering
\small
\begin{verbatim}
Telegram API
     │
     ▼
┌──────────────────────────────────────────────────┐
│  МАСТЕР  (мастер/cl/)                            │
│  SBCL-процесс #1                                 │
│  - poll-loop: long-polling Telegram              │
│  - роутинг команд: /породить /запустить ...      │
│  - душа, история диалогов, комбинирование        │
└──────────────────────────────────────────────────┘
         │ load .lisp
         ▼
┌──────────────────────────────────────────────────┐
│  мастер/потоки/                                  │
│  *.lisp — файлы потоков                          │
│  порождатель.lisp  эхо.lisp  уточнение.lisp      │
│  питон.lisp        реакт.lisp                    │
└──────────────────────────────────────────────────┘
         │ load .lisp (общий каталог)
         ▼
┌──────────────────────────────────────────────────┐
│  ПОДМАСТЕРЬЕ  (подмастерье/)                     │
│  SBCL-процесс #2                                 │
│  - персистентный образ                           │
│  - загружает потоки при первом обращении         │
│  - вызывает (выполнить задача) напрямую          │
└──────────────────────────────────────────────────┘
\end{verbatim}
\caption{Архитектура Development System}
\label{fig:arch}
\end{figure}

\subsection{Файловая структура}

\begin{table}[H]
\centering\small
\begin{tabular}{lll}
\toprule
\textbf{Путь} & \textbf{Язык} & \textbf{Роль} \\
\midrule
\texttt{мастер/cl/мастер.asd}     & ASDF  & Определение системы \\
\texttt{мастер/cl/packages.lisp}  & CL    & Пакет \texttt{:мастер} \\
\texttt{мастер/cl/config.lisp}    & CL    & Переменные окружения \\
\texttt{мастер/cl/telegram.lisp}  & CL    & HTTP-клиент Telegram \\
\texttt{мастер/cl/llm.lisp}       & CL    & HTTP-клиент OpenRouter \\
\texttt{мастер/cl/dusha.lisp}     & CL    & Личность бота (soul) \\
\texttt{мастер/cl/история.lisp}   & CL    & История диалогов \\
\texttt{мастер/cl/потоки.lisp}    & CL    & Управление потоками \\
\texttt{мастер/cl/комбо.lisp}     & CL    & Комбинирование потоков \\
\texttt{мастер/cl/команды.lisp}   & CL    & Обработчики команд \\
\texttt{мастер/cl/main.lisp}      & CL    & Poll-loop, точка входа \\
\texttt{подмастерье/main.lisp}    & CL    & Персистентный исполнитель \\
\texttt{подмастерье/исполнитель.lisp} & CL & Загрузка и запуск потоков \\
\texttt{подмастерье/tracer.lisp}  & CL    & Трассировка вычислений \\
\texttt{мастер/потоки/*.lisp}     & CL    & Вычислительные потоки \\
\bottomrule
\end{tabular}
\caption{Файловая структура Development System}
\end{table}

\subsection{Контракт потока}

Каждый поток — это обычный CL-файл, соответствующий одному контракту:

\begin{lispcode}[caption={Контракт вычислительного потока}]
(defpackage :поток-ИМЯ
  (:use :cl)
  (:export #:выполнить))

(in-package :поток-ИМЯ)

(defun выполнить (задача)
  ;; задача — строка
  ;; возвращает строку-результат
  ...)
\end{lispcode}

Соглашение предельно простое: один пакет, одна экспортируемая функция. Это позволяет как Мастеру, так и Подмастерью находить точку входа через \texttt{find-symbol}:

\begin{lispcode}[caption={Поиск точки входа через рефлексию}]
(defun %найти-выполнить (имя)
  (let ((pkg (find-package (string-upcase имя))))
    (when pkg
      (let ((sym (find-symbol "ВЫПОЛНИТЬ" pkg)))
        (when (and sym (fboundp sym))
          (symbol-function sym))))))
\end{lispcode}

% ═══════════════════════════════════════════════════════════════════════════════
\section{Мастер: Telegram-оркестратор}
\label{sec:master}

\subsection{ASDF-система и пакет}

Мастер оформлен как полноценная ASDF-система с явным списком зависимостей и порядком загрузки файлов:

\begin{lispcode}[caption={мастер/cl/мастер.asd}]
(asdf:defsystem #:мастер
  :description "Telegram bot / orchestrator"
  :depends-on (#:dexador #:cl-json)
  :serial t
  :components ((:file "packages")
               (:file "config")
               (:file "telegram")
               (:file "llm")
               (:file "dusha")
               (:file "история")
               (:file "потоки")
               (:file "комбо")
               (:file "команды")
               (:file "main")))
\end{lispcode}

Ключ \texttt{:serial t} гарантирует, что файлы компилируются в порядке объявления. Это важно: \texttt{команды.lisp} вызывает функции из \texttt{потоки.lisp} и \texttt{история.lisp}, которые должны быть скомпилированы раньше.

Пакет \texttt{:мастер} экспортирует все публичные точки входа:

\begin{lispcode}[caption={мастер/cl/packages.lisp}]
(defpackage #:мастер
  (:use #:cl #:uiop)
  (:nicknames #:м)
  (:export
   #:start #:poll-loop
   #:обработать-команду #:обработать-update
   #:читать-душу #:писать-душу #:душа->системный-промпт
   #:llm-complete
   #:get-updates #:send-message #:send-document
   #:answer-callback-query #:make-inline-keyboard
   #:загрузить-историю #:сохранить-историю
   #:добавить-сообщение #:очистить-историю
   #:список-потоков #:загрузить-поток
   #:загрузить-все-потоки #:запустить-поток
   #:активные-потоки #:комбинировать))
\end{lispcode}

\subsection{Telegram HTTP-клиент}

Telegram Bot API работает через HTTPS с JSON-телом запроса. Мастер использует \texttt{dexador} — современную CL-библиотеку для HTTP:

\begin{lispcode}[caption={мастер/cl/telegram.lisp — ядро HTTP-клиента}]
(defun api-url (method)
  (concatenate 'string *telegram-api-base* *токен* "/" method))

(defun api-post (method payload)
  "POST JSON payload к Telegram API.
  Возвращает распарсенный alist или nil."
  (handler-case
      (let* ((body (cl-json:encode-json-to-string payload))
             (resp (dexador:post (api-url method)
                                 :content body
                                 :headers '(("Content-Type"
                                             . "application/json")))))
        (cl-json:decode-json-from-string resp))
    (error (e)
      (format *error-output* "[telegram] ошибка ~A: ~A~%" method e)
      nil)))
\end{lispcode}

Примечательно: вся обработка ошибок сводится к одному \texttt{handler-case}. При любом сбое (таймаут, 4xx, 5xx) функция возвращает \texttt{nil} вместо сигнала условия. Это осознанный выбор: poll-loop должен быть устойчивым к сетевым флуктуациям.

Long-polling реализован через параметр \texttt{timeout=30} в \texttt{getUpdates}:

\begin{lispcode}[caption={мастер/cl/telegram.lisp — long-poll}]
(defun get-updates (&key (timeout 30) offset)
  "Long-poll Telegram. Возвращает список update-alist или nil."
  (let* ((params (append `((:timeout . ,timeout))
                         (when offset `((:offset . ,offset)))))
         (resp (api-post "getUpdates" params)))
    (when resp (cdr (assoc :result resp)))))
\end{lispcode}

Telegram удерживает соединение до 30 секунд, если нет новых сообщений, затем возвращает пустой массив. Это значительно снижает количество запросов по сравнению с short-polling.

\subsection{Главный цикл}

Poll-loop реализован через хвостовую рекурсию. В Common Lisp хвостовая рекурсия не гарантируется стандартом, но SBCL оптимизирует её при включённой скорости:

\begin{lispcode}[caption={мастер/cl/main.lisp — рекурсивный poll-loop}]
(defun poll-loop (offset)
  (when *polling*
    (handler-case
        (let* ((updates (get-updates :offset offset :timeout 30))
               (max-id offset))
          (dolist (upd (or updates '()))
            (let ((id (cdr (assoc :update--id upd))))
              (when (and id (> id max-id))
                (setf max-id id)))
            (ignore-errors (обработать-update upd)))
          (poll-loop
            (if (> max-id offset) (1+ max-id) offset)))
      (error (e)
        (format *error-output*
                "[мастер] ошибка в poll-loop: ~A~%" e)
        (sleep 2)
        (poll-loop offset)))))
\end{lispcode}

Offset-управление следует протоколу Telegram: следующий запрос отправляется с \texttt{offset = max\_update\_id + 1}, что подтверждает получение всех предыдущих сообщений.

Обработка входящих update-ов диспетчеризируется по типу: сообщение (\texttt{:message}) или callback от inline-кнопки (\texttt{:callback--query}):

\begin{lispcode}[caption={мастер/cl/main.lisp — dispatch по типу update}]
(defun обработать-update (update)
  (cond
    ((assoc :message update)
     (let* ((msg     (cdr (assoc :message update)))
            (chat    (cdr (assoc :chat msg)))
            (chat-id (cdr (assoc :id chat)))
            (text    (cdr (assoc :text msg))))
       (when (stringp text)
         (send-message chat-id
           (обработать-команду chat-id text)))))
    ((assoc :callback--query update)
     (let* ((cb      (cdr (assoc :callback--query update)))
            (chat-id (cdr (assoc :id
                            (cdr (assoc :chat
                              (cdr (assoc :message cb)))))))
            (data    (cdr (assoc :data cb))))
       (answer-callback-query (cdr (assoc :id cb)))
       (send-message chat-id
         (обработать-команду chat-id data))))))
\end{lispcode}

Заметим двойное дефисирование в ключах (\texttt{:update--id}, \texttt{:callback--query}): это особенность \texttt{cl-json}, который конвертирует \texttt{snake\_case} JSON-ключи в символы CL, заменяя подчёркивания двойными дефисами.

\subsection{Роутинг команд}

Команды хранятся в \texttt{assoc}-таблице: символ команды → функция-обработчик. Это классический паттерн «dispatch table»:

\begin{lispcode}[caption={мастер/cl/команды.lisp — dispatch table}]
(defparameter *команды*
  (list (cons "/старт"       #'%обр-старт)
        (cons "/start"       #'%обр-старт)
        (cons "/потоки"      #'%обр-потоки)
        (cons "/запустить"   #'%обр-запустить)
        (cons "/остановить"  #'%обр-остановить)
        (cons "/переключить" #'%обр-переключить)
        (cons "/состояние"   #'%обр-состояние)
        (cons "/породить"    #'%обр-породить)
        (cons "/диалог"      #'%обр-диалог)
        (cons "/сбросить"    #'%обр-сбросить)))

(defun обработать-команду (chat-id text)
  (multiple-value-bind (cmd args) (%cmd-args text)
    (let ((fn (cdr (assoc cmd *команды* :test #'string=))))
      (if fn
          (handler-case (funcall fn chat-id args)
            (error (e) (format nil "Ошибка: ~A" e)))
          (format nil "Неизвестная команда: ~A" cmd)))))
\end{lispcode}

Парсинг строки команды выполняется самостоятельно, без внешних библиотек:

\begin{lispcode}[caption={мастер/cl/команды.lisp — разбор строки команды}]
(defun %split-words (str)
  (let ((result '())
        (cur (make-array 0 :element-type 'character
                           :adjustable t :fill-pointer 0)))
    (loop for ch across str do
      (if (char= ch #\Space)
          (when (plusp (length cur))
            (push (copy-seq cur) result)
            (setf (fill-pointer cur) 0))
          (vector-push-extend ch cur)))
    (when (plusp (length cur))
      (push (copy-seq cur) result))
    (nreverse result)))
\end{lispcode}

\subsection{Управление состоянием: душа, история, потоки}

\subsubsection{Душа (Soul)}

«Душа» — это S-выражение в файле \texttt{душа.scm}, определяющее личность бота: имя, описание, стиль, инструкции. Файл читается стандартным \texttt{read}:

\begin{lispcode}[caption={мастер/cl/dusha.lisp}]
(defun читать-душу (путь)
  (with-open-file (in путь :direction :input)
    (read in)))

(defun писать-душу (душа путь)
  (with-open-file (out путь :direction :output
                   :if-exists :supersede
                   :if-does-not-exist :create)
    (let ((*print-pretty* t))
      (prin1 душа out))))

(defun душа->системный-промпт (душа)
  (let ((имя         (getf душа :имя))
        (описание    (getf душа :описание))
        (стиль       (getf душа :стиль))
        (инструкции  (getf душа :инструкции)))
    (format nil "Имя: ~a~%Описание: ~a~%Стиль: ~a~%~
                 Инструкции: ~a"
            имя описание стиль инструкции)))
\end{lispcode}

Формат \texttt{душа.scm} — это plist:
\begin{verbatim}
(:имя "Мастер" :описание "Оркестратор потоков"
 :стиль "лаконичный" :инструкции "Отвечай кратко и по делу.")
\end{verbatim}

\subsubsection{История диалогов}

История хранится в \texttt{история/<chat-id>.json} и представляет собой список сообщений в формате OpenAI (\texttt{role}/\texttt{content}):

\begin{lispcode}[caption={мастер/cl/история.lisp}]
(defun %путь-истории (chat-id)
  (merge-pathnames (format nil "~A.json" chat-id)
                   *каталог-истории*))

(defun добавить-сообщение (chat-id role content)
  (let* ((история (or (загрузить-историю chat-id) '()))
         (новая (append история
                  (list `((:role . ,role)
                          (:content . ,content))))))
    (сохранить-историю chat-id новая)
    новая))

(defun загрузить-историю (chat-id)
  (let ((файл (%путь-истории chat-id)))
    (when (probe-file файл)
      (handler-case
          (cl-json:decode-json-from-string
           (uiop:read-file-string файл))
        (error () nil)))))
\end{lispcode}

\subsubsection{Управление потоками}

Состояние активных потоков хранится в \texttt{hash-table}. Поиск функции \texttt{выполнить} реализован через рефлексию — \texttt{find-package} и \texttt{find-symbol}:

\begin{lispcode}[caption={мастер/cl/потоки.lisp}]
(defparameter *активные-потоки*
  (make-hash-table :test #'equal))

(defun список-потоков ()
  (mapcar #'pathname-name
    (directory (merge-pathnames "*.lisp" *каталог-потоков*))))

(defun %найти-выполнить (имя)
  (let ((pkg (find-package (string-upcase имя))))
    (when pkg
      (let ((sym (find-symbol "ВЫПОЛНИТЬ" pkg)))
        (when (and sym (fboundp sym))
          (symbol-function sym))))))

(defun запустить-поток (имя задача)
  (let* ((путь (merge-pathnames (format nil "~A.lisp" имя)
                                *каталог-потоков*)))
    (when (probe-file путь)
      (загрузить-поток путь)
      (let ((fn (%найти-выполнить имя)))
        (when fn
          (setf (gethash имя *активные-потоки*) t)
          (funcall fn задача))))))
\end{lispcode}

\subsubsection{Комбинирование потоков}

Функция \texttt{комбинировать} поддерживает два режима: последовательный (pipeline) и параллельный:

\begin{lispcode}[caption={мастер/cl/комбо.lisp}]
(defun %join (sep strings)
  (reduce (lambda (a b) (concatenate 'string a sep b))
          strings :initial-value ""))

(defun комбинировать (потоки задача &key параллельно)
  (if параллельно
      (%join "\n---\n"
             (mapcar (lambda (имя)
                       (or (ignore-errors
                             (запустить-поток имя задача))
                           ""))
                     потоки))
      (reduce (lambda (вход имя)
                (or (ignore-errors
                      (запустить-поток имя вход))
                    вход))
              потоки :initial-value задача)))
\end{lispcode}

В последовательном режиме \texttt{reduce} элегантно реализует pipeline: результат каждого потока становится входом следующего.

\subsection{LLM-клиент}

Запросы к языковой модели через OpenRouter:

\begin{lispcode}[caption={мастер/cl/llm.lisp}]
(defun llm-complete (prompt &key (model *openrouter-model*)
                               system messages)
  (let* ((msgs (append
                 (when system
                   (list `(("role" . "system")
                           ("content" . ,system))))
                 (if messages messages
                     (list `(("role" . "user")
                             ("content" . ,prompt))))))
         (body (cl-json:encode-json-to-string
                 `(("model" . ,model) ("messages" . ,msgs))))
         (headers `(("Authorization" . ,(format nil "Bearer ~a"
                                         *openrouter-key*))
                    ("Content-Type" . "application/json")))
         (resp (dexador:post *openrouter-api-url*
                :content body :headers headers))
         (json (ignore-errors
                 (cl-json:decode-json-from-string resp)))
         (choices (cdr (assoc :choices json))))
    (if (and choices (consp choices))
        (let ((msg (cdr (assoc :message (first choices)))))
          (or (cdr (assoc :content msg)) "[нет ответа]"))
        "[нет ответа]")))
\end{lispcode}

% ═══════════════════════════════════════════════════════════════════════════════
\section{Подмастерье: персистентный исполнитель}
\label{sec:apprentice}

\subsection{Концепция персистентного образа}

Традиционный подход к выполнению кода по запросу — запуск нового процесса на каждый вызов. Это удобно с точки зрения изоляции, но дорого: SBCL стартует около 400ms, плюс время на загрузку зависимостей.

Подмастерье использует обратный подход: \textbf{один долгоживущий процесс, в который потоки загружаются динамически}. После первой загрузки поток живёт в образе и вызывается напрямую через \texttt{funcall}. Повторный вызов — это просто вызов функции, без файлового ввода-вывода.

\subsection{Явное состояние (SICP-style)}

Состояние Подмастерья хранится как явный параметр, передаваемый через рекурсию. Это паттерн из SICP (Structure and Interpretation of Computer Programs) — мутабельное состояние заменяется аккумулятором в хвостовой рекурсии:

\begin{lispcode}[caption={подмастерье/main.lisp — рекурсивный цикл с явным состоянием}]
(defun опрос (сдвиг реестр)
  (handler-case
      (let* ((тело       (получить-обновления сдвиг))
             (обновления (or (разобрать-обновления тело) '()))
             (сообщения  (извлечь-сообщения обновления)))
        (multiple-value-bind (новый-рег новый-сдвиг)
            (обр-сообщения сообщения реестр сдвиг)
          (when (> новый-сдвиг сдвиг)
            (сохранить-сдвиг новый-сдвиг))
          (опрос новый-сдвиг новый-рег)))
    (error (e)
      (format t "Ошибка цикла: ~a~%" e)
      (sleep 5)
      (опрос сдвиг реестр))))
\end{lispcode}

Здесь \texttt{реестр} — это alist вида \texttt{(имя-пакета . путь)}, который накапливается по мере загрузки новых потоков. Обратим внимание: глобального мутабельного состояния нет. Каждая итерация получает реестр как параметр и возвращает обновлённый.

\subsection{Реестр потоков и активный поток}

Активный поток определяется через файл \texttt{активный.scm} — S-выражение, которое Мастер записывает при переключении потока:

\begin{lispcode}[caption={подмастерье/main.lisp — чтение активного потока}]
(defun прочитать-активный ()
  (let ((путь (format nil "~aактивный.scm" *каталог-потоков*)))
    (handler-case
        (with-open-file (f путь)
          (let* ((д     (read f nil nil))
                 (файл  (cadr (assoc 'файл  (cdr д))))
                 (пакет (cadr (assoc 'пакет (cdr д)))))
            (when (and файл пакет)
              (list файл пакет))))
      (error () nil))))
\end{lispcode}

Формат \texttt{активный.scm}:
\begin{verbatim}
(активный
  (файл  "/path/to/потоки/уточнение.lisp")
  (пакет "уточнение"))
\end{verbatim}

\subsection{Загрузка потока в образ}

Загрузка потока идемпотентна — если поток уже в реестре, он не перегружается:

\begin{lispcode}[caption={подмастерье/main.lisp — ленивая загрузка в образ}]
(defun загрузить-в-образ (путь пакет реестр)
  "Загружает поток в образ если ещё не в реестре.
  Возвращает новый реестр."
  (if (assoc пакет реестр :test #'string=)
      реестр
      (handler-case
          (progn
            (load путь :verbose nil :print nil)
            (format t "Загружен поток: ~a~%" пакет)
            (cons (cons пакет путь) реестр))
        (error (e)
          (format t "Ошибка загрузки ~a: ~a~%" пакет e)
          реестр))))
\end{lispcode}

\subsection{Модуль исполнителя}

Файл \texttt{исполнитель.lisp} предоставляет обобщённый API для загрузки и запуска потоков с развитой обработкой условий. Система рестартов позволяет восстанавливаться от ошибок без перезапуска процесса:

\begin{lispcode}[caption={подмастерье/исполнитель.lisp — выполнение с рестартами}]
(defun выполнить-задачу (поток-функция задача)
  "Вызывает функцию потока с поддержкой рестартов."
  (restart-case
      (handler-case
          (funcall поток-функция задача)
        (error (у)
          (format nil "Ошибка выполнения: ~a" у)))
    (повторить ()
      :report "Повторить выполнение задачи"
      (выполнить-задачу поток-функция задача))
    (подменить (значение)
      :report "Использовать запасное значение"
      :interactive (lambda ()
                     (format t "Введите запасное значение: ")
                     (list (read)))
      значение)))
\end{lispcode}

Функция \texttt{выполнить-с-повтором} реализует retry-логику через хвостовую рекурсию — без \texttt{loop}, без мутабельного счётчика:

\begin{lispcode}[caption={подмастерье/исполнитель.lisp — retry через рекурсию}]
(defun выполнить-с-повтором (поток-функция задача макс-попыток)
  (let ((попытка 0))
    (labels ((попробовать ()
               (incf попытка)
               (handler-case
                   (funcall поток-функция задача)
                 (error (у)
                   (if (< попытка макс-попыток)
                       (попробовать)
                       (format nil "Ошибка после ~a попыток: ~a"
                               макс-попыток у))))))
      (попробовать))))
\end{lispcode}

% ═══════════════════════════════════════════════════════════════════════════════
\section{Поток-порождатель: метациркулярная петля}
\label{sec:spawner}

\subsection{Идея}

Порождатель — это поток, который умеет создавать новые потоки. Он принимает описание на естественном языке, отправляет его LLM с системным промптом-шаблоном, получает CL-код, компилирует его и загружает в тот же образ. Тем самым петля замыкается: \textit{исполняемая среда порождает сама себя}.

\subsection{Системный промпт как контракт}

Промпт задаёт жёсткий контракт на структуру генерируемого кода:

\begin{lispcode}[caption={мастер/потоки/порождатель.lisp — системный промпт}]
(defparameter *промпт-системы*
  "Ты генератор вычислительных потоков на Common Lisp.
Верни ТОЛЬКО код Common Lisp без пояснений, без markdown.
Требования к потоку:
- (defpackage :поток-ИМЯ (:use :cl) (:export #:выполнить))
- (in-package :поток-ИМЯ) после defpackage
- (defun выполнить (задача) ...) — единственная точка входа
- OPENROUTER_API_KEY через (sb-ext:posix-getenv ...)
- HTTP через drakma, JSON через cl-json (уже загружены)
- рекурсия вместо loop, функции до 15 строк")
\end{lispcode}

\subsection{Извлечение и компиляция кода}

LLM иногда оборачивает код в markdown-блок (\texttt{```lisp...```}). Для обработки этого случая реализован парсер через рекурсию:

\begin{lispcode}[caption={мастер/потоки/порождатель.lisp — извлечение кода из LLM-ответа}]
(defun начало-блока-p (строка)
  (let ((с (string-trim '(#\Space #\Tab) строка)))
    (and (>= (length с) 3)
         (string= (subseq с 0 3) "```"))))

(defun извлечь-из-блока (строки)
  "Убирает ```-обёртку если есть."
  (labels ((ищи (ост внутри накоп)
             (cond
               ((null ост)
                (when накоп
                  (format nil "~{~a~%~}" (reverse накоп))))
               ((and (not внутри) (начало-блока-p (car ост)))
                (ищи (cdr ост) t накоп))
               ((and внутри (начало-блока-p (car ост)))
                (ищи (cdr ост) nil накоп))
               (внутри
                (ищи (cdr ост) t (cons (car ост) накоп)))
               (t
                (ищи (cdr ост) nil накоп)))))
    (let ((строки-список
           (uiop:split-string строки :separator '(#\Newline))))
      (or (ищи строки-список nil nil) строки))))
\end{lispcode}

Имя пакета извлекается из кода через поиск подстроки \texttt{"поток-"}:

\begin{lispcode}[caption={мастер/потоки/порождатель.lisp — извлечение имени пакета}]
(defun имя-пакета (код)
  (let* ((метка "поток-")
         (поз   (search метка код)))
    (if (null поз)
        "безымянный"
        (let* ((нач (+ поз (length метка)))
               (кон (or (position-if
                          (lambda (c)
                            (member c '(#\Space #\Newline
                                        #\( #\) #\` #\")))
                          код :start нач)
                        (length код))))
          (subseq код нач кон)))))
\end{lispcode}

\subsection{Загрузка в живой образ}

После успешной компиляции поток загружается в тот же образ, где работает Мастер:

\begin{lispcode}[caption={мастер/потоки/порождатель.lisp — загрузка результата в образ}]
(defun выполнить (описание)
  (handler-case
      (let* ((сырой (запросить-генерацию описание))
             (код   (извлечь-из-блока сырой))
             (имя   (имя-пакета код))
             (путь  (записать-поток имя код))
             (ok?   (компилировать путь))
             (знак  (if ok? "OK" "WARN")))
        (when ok?
          (загрузить-в-образ путь
            (format nil "поток-~a" имя)))
        (format nil "Поток ~s ~a ~a и доступен в образе."
                имя знак (if ok? "скомпилирован" "не скомпилирован")))
    (error (e)
      (format nil "Ошибка порождения: ~a" e))))
\end{lispcode}

Функция \texttt{компилировать} использует \texttt{compile-file} и проверяет третье возвращаемое значение (\texttt{failure-p}):

\begin{lispcode}[caption={мастер/потоки/порождатель.lisp — компиляция}]
(defun компилировать (путь)
  (handler-case
      (multiple-value-bind (_ warn fail)
          (compile-file путь :print nil :verbose nil)
        (declare (ignore _ warn))
        (not fail))
    (error () nil)))
\end{lispcode}

% ═══════════════════════════════════════════════════════════════════════════════
\section{Трассировка вычислений}
\label{sec:tracer}

\texttt{tracer.lisp} реализует механизм записи хода вычислений в файл \texttt{.след} как серии S-выражений. Это позволяет анализировать пошаговую работу потока post-hoc.

\subsection{API трассировки}

\begin{lispcode}[caption={подмастерье/tracer.lisp — API трассировки}]
(defpackage :трас
  (:use :cl)
  (:export :начать-след :записать-шаг :завершить-след
           :*путь-журнала* :*имя-потока* :*задача*))

(defvar *путь-журнала* nil "Путь к файлу журнала.")
(defvar *шаги*         nil "Накопленные шаги.")
(defvar *начало*       0   "Время начала в мс.")

(defun записать-шаг (номер вход ответ &key (качество :неизвестно))
  (let ((мс (- (мс-сейчас) *начало*)))
    (push `(шаг (номер ,номер)
                (вход  ,вход)
                (ответ ,ответ)
                (мс    ,мс)
                (качество ,качество))
          *шаги*)))
\end{lispcode}

\subsection{Формат .след файлов}

Итоговый журнал — читаемое S-выражение, которое можно загрузить в REPL через \texttt{read}:

\begin{lispcode}[caption={Пример .след файла}]
(след
  (поток "уточнение") (время "2026-03-13T11:30:00")
  (задача "разберись что такое continuation")
  (шаги
    (шаг (номер 1) (вход "разберись что такое continuation")
         (ответ "continuation — это...") (мс 1234)
         (качество :недостаточно))
    (шаг (номер 2) (вход "уточни про call/cc")
         (ответ "call/cc в Scheme...") (мс 2100)
         (качество :достаточно)))
  (итог "Объяснение continuation дано.") (мс-всего 3334)
  (итераций 2))
\end{lispcode}

\subsection{Инициализация трассировщика}

Трассировщик загружается Подмастерьем автоматически при наличии переменной окружения \texttt{TRACE\_FILE}:

\begin{lispcode}[caption={подмастерье/исполнитель.lisp — ленивая загрузка трассировщика}]
(defun обеспечить-трассировщик ()
  (let* ((путь-трас
          (merge-pathnames "tracer.lisp"
            (make-pathname
              :defaults *load-pathname*
              :name nil :type nil)))
         (путь-журнала (sb-ext:posix-getenv "TRACE_FILE")))
    (when (and путь-журнала (probe-file путь-трас))
      (load путь-трас :verbose nil :print nil)
      (when (find-package :трас)
        (setf (symbol-value
                (find-symbol "*ПУТЬ-ЖУРНАЛА*" :трас))
              путь-журнала)))))
\end{lispcode}

% ═══════════════════════════════════════════════════════════════════════════════
\section{Ключевые архитектурные решения}
\label{sec:design}

\subsection{Единый язык в стеке}

Переход с Racket на Common Lisp устранил несоответствие языков: оркестратор и потоки теперь в одном языковом пространстве. Это открывает возможности, недостижимые в полиязычной архитектуре:

\begin{itemize}[noitemsep]
  \item Поток, загруженный в образ Мастера, напрямую вызывается через \texttt{funcall} — без сериализации, без IPC;
  \item Общие библиотеки (\texttt{drakma}, \texttt{cl-json}) загружаются один раз;
  \item Ошибки в потоках ловятся системой условий CL — не падают весь процесс.
\end{itemize}

\subsection{Контракт вместо фреймворка}

Вместо тяжёлого фреймворка система использует минимальный контракт: \texttt{defpackage :поток-ИМЯ} + \texttt{defun выполнить (задача)}. Это:
\begin{itemize}[noitemsep]
  \item Делает потоки независимыми файлами, которые можно тестировать в REPL;
  \item Позволяет LLM генерировать потоки без знания внутренностей системы;
  \item Исключает coupling: Подмастерье не знает ничего о содержимом потока.
\end{itemize}

\subsection{Ленивая загрузка и идемпотентность}

Поток загружается в образ только при первом обращении и помещается в реестр. Повторные вызовы работают с уже загруженной функцией. Это инверсия типичного serverless-паттерна: \textit{не «запустить процесс на запрос», а «держать всё горячим»}.

\subsection{Явное состояние vs. мутабельные переменные}

Цикл опроса Подмастерья передаёт состояние (реестр потоков, сдвиг) как параметры рекурсии. Это устраняет целый класс ошибок: невозможно случайно мутировать реестр из другого потока исполнения, невозможно забыть сбросить сдвиг при рестарте.

Исключение: \texttt{*активные-потоки*} в Мастере — это \texttt{hash-table}, потому что его модифицируют разные обработчики команд, и передавать его через рекурсию poll-loop означало бы усложнить сигнатуру без выигрыша.

\subsection{Метациркулярность как проектное решение}

Поток-порождатель замыкает петлю: система способна расширять саму себя в рантайме без перезапуска и без вмешательства разработчика. Это не хак, а проектное решение: Common Lisp — гомоиконный язык, где код и данные имеют одно представление. \texttt{compile-file} + \texttt{load} — штатные операции над образом.

% ═══════════════════════════════════════════════════════════════════════════════
\section{Запуск и развёртывание}

\subsection{Зависимости}

\begin{itemize}[noitemsep]
  \item SBCL ≥ 2.0
  \item Quicklisp (загружается автоматически из \texttt{\textasciitilde/quicklisp/setup.lisp})
  \item \texttt{dexador} — HTTP-клиент (через Quicklisp)
  \item \texttt{cl-json} — JSON (через Quicklisp)
  \item \texttt{drakma} + \texttt{flexi-streams} — HTTP для Подмастерья (через Quicklisp)
\end{itemize}

\subsection{Переменные окружения}

\begin{table}[H]
\centering\small
\begin{tabular}{lll}
\toprule
\textbf{Переменная} & \textbf{Компонент} & \textbf{Назначение} \\
\midrule
\texttt{MASTER\_BOT\_TOKEN}     & Мастер       & Telegram Bot API токен \\
\texttt{APPRENTICE\_BOT\_TOKEN} & Подмастерье  & Telegram Bot API токен \\
\texttt{ADMIN\_CHAT\_ID}        & Оба          & Whitelist chat ID \\
\texttt{OPENROUTER\_API\_KEY}   & Оба          & Ключ к OpenRouter \\
\texttt{TRACE\_FILE}            & Подмастерье  & Путь к файлу трассировки \\
\bottomrule
\end{tabular}
\caption{Переменные окружения}
\end{table}

\subsection{Скрипт запуска Мастера}

\begin{lispcode}[language=bash,caption={мастер/cl/start.sh}]
#!/usr/bin/env bash
export MASTER_BOT_TOKEN="${MASTER_BOT_TOKEN:?нужен MASTER_BOT_TOKEN}"
export ADMIN_CHAT_ID="${ADMIN_CHAT_ID:?нужен ADMIN_CHAT_ID}"
export OPENROUTER_API_KEY="${OPENROUTER_API_KEY:?нужен OPENROUTER_API_KEY}"

SCRIPT_DIR="$(cd "$(dirname "$0")" && pwd)"

exec sbcl --noinform \
  --load ~/quicklisp/setup.lisp \
  --eval "(push #p\"${SCRIPT_DIR}/\" asdf:*central-registry*)" \
  --eval '(asdf:load-system :мастер :verbose nil)' \
  --eval '(мастер:start)'
\end{lispcode}

% ═══════════════════════════════════════════════════════════════════════════════
\section{Тестирование}

Система покрыта набором smoke-тестов, проверяющих корректность ключевых функций без запуска реального бота:

\begin{lispcode}[caption={мастер/cl/тесты.lisp — smoke-тесты}]
;; 1. api-url формат
(тест "api-url формат"
  (let ((url (api-url "getUpdates")))
    (assert (search "telegram.org/bot" url))
    (assert (search "getUpdates" url))))

;; 2. душа->системный-промпт
(тест "душа->системный-промпт"
  (let* ((ф '(:имя "Тест" :описание "тест-бот"
              :стиль "краткий" :инструкции "коротко"))
         (п (душа->системный-промпт ф)))
    (assert (search "Тест" п))
    (assert (search "краткий" п))))

;; 3. список-потоков
(тест "список-потоков возвращает строки"
  (let ((lst (список-потоков)))
    (assert (listp lst))
    (dolist (имя lst) (assert (stringp имя)))))

;; 4. роутинг команды
(тест "обработать-команду /старт"
  (let ((resp (обработать-команду 0 "/старт")))
    (assert (stringp resp))
    (assert (plusp (length resp)))))
\end{lispcode}

Результат прогона:
\begin{verbatim}
[ТЕСТЫ МАСТЕР]
  ✓ api-url формат
  ✓ душа->системный-промпт
  ✓ список-потоков возвращает строки
  ✓ обработать-команду /старт
  ✓ %split-words базовый
  ✓ обработать-команду неизвестная → строка
Итого ошибок: 0
\end{verbatim}

% ═══════════════════════════════════════════════════════════════════════════════
\section{Заключение}
\label{sec:conclusion}

Статья описала Development System — двухкомпонентную систему оркестрации LLM-потоков на Common Lisp. Основные результаты:

\begin{enumerate}[noitemsep]
  \item Сформулирован и реализован \textbf{контракт потока}: \texttt{defpackage :поток-ИМЯ} + \texttt{defun выполнить (задача) → строка}. Контракт прост, проверяем и достаточен для полной автоматизации запуска.

  \item Реализован \textbf{персистентный образ} с ленивой загрузкой: потоки загружаются в SBCL-образ один раз и вызываются напрямую без subprocess-оверхеда.

  \item Реализована \textbf{метациркулярная петля}: поток-порождатель генерирует новые потоки через LLM, компилирует их и загружает в тот же образ, не требуя перезапуска процесса.

  \item Система \textbf{полностью мигрирована с Racket} на Common Lisp. Единый язык в стеке устранил несоответствие семантик и дал доступ к нативной системе условий SBCL.

  \item Архитектура следует принципу \textbf{явного состояния} (SICP-style): цикл опроса передаёт реестр потоков как параметр рекурсии, исключая скрытые глобальные мутации.
\end{enumerate}

Направления дальнейшего развития: интеграция формальной верификации через ACL2 (доказательство замкнутости потоков), реализация протокола \texttt{ask-user} для интерактивных checkpoint-ов в середине выполнения, и расширение трассировщика для генерации человекочитаемых отчётов.

\end{document}
